\section{Exact division and integer square root}
\label{sec:exact-arithmetic-proof}

Let $n$ be the supplied input width, let $a$ be the nonnegative dividend
or radicand, and let $b>0$ be the divisor for the division task, with
the values and output conventions of Corollary~\ref{cor:exact-arithmetic}.
Here $p\ge1$ denotes an integer target precision for a Newton routine.

\begin{proof}[Proof of Corollary~\ref{cor:exact-arithmetic}]
We use the classical Newton reductions as presented by
Brent~\cite[Lemmas 2.1--2.4]{Brent1976}, with their tape costs made
explicit. Set $B(p)=p(\lg p)^{1-\kappa}$. A linear scan handles a zero
dividend or radicand and finds the significant bits of a positive input.
For division write $b=2^h c$, where $0\le h<n$ and $1\le c<2$.
For square root write $a=2^{2e}c$, where $0\le e<n/2$ and $1\le c<4$.
Copy the normalized digits most significant first to a tape with a marked
beginning. This initial scan also records the relevant exponent.

For an integer $t\ge1$, let
$R_t(z)=2^{-t}\lfloor 2^t z\rfloor$ and put $\delta=2^{-t}$.
The prefix $c_t=R_t(c)$ has $0\le c-c_t<\delta$.
Each product below is computed exactly on integer significands and
then rounded by the indicated $R_t$; additions and shifts are exact
until a stated rounding. When precision increases, the stored dyadic
iterate is extended by zero fractional bits.

For the reciprocal iteration $x\leftarrow x(2-cx)$, compute
\[
 u=R_t(c_t x),\qquad x^+=R_t\bigl(x(2-u)\bigr).
\]
Put $\varepsilon=1-cx$. If $|\varepsilon|\le1/2$, then
$0<x\le3/2$. Writing $\eta_1=u-c_t x$ and
$\eta_2=x^+-x(2-u)$ gives $|\eta_i|<\delta$ and
\[
 1-cx^+
 =\varepsilon^2-c(c-c_t)x^2+cx\eta_1-c\eta_2.
\]
Since $cx^2=(1-\varepsilon)^2/c\le9/4$, $cx\le3/2$, and $c<2$,
\[
 |1-cx^+|
 \le\varepsilon^2+\frac{23}{4}\delta
 <\varepsilon^2+6\delta.                                  \tag{*}
\]
The fixed seed $x=5/8$ has $|\varepsilon|\le3/8$ throughout
$1\le c<2$. One step with $t=6$ therefore gives residual less than
$9/64+6/64=15/64<1/4$.

For a requested target $p\ge1$, begin with $m=1$ after this warm-up.
While $m<p$, set $m'=\min\{2m,p\}$ and work at $t=m'+5$ fractional
bits. If the residual is at most $2^{-m-1}$, then $(*)$ gives
\[
 |1-cx^+|
 \le2^{-2m-2}+6\,2^{-m'-5}
 \le\frac7{16}\,2^{-m'}<2^{-m'-1}.
\]
Set $m=m'$ and continue. The reciprocal routine at target $p$ thus
returns a dyadic $x$ satisfying
$|1-cx|\le2^{-p-1}$ and $|x-1/c|\le2^{-p-1}$.
These errors are measured against the true $c$ at every stage, although
the prefix $c_t$ changes.

For the root iteration $x\leftarrow(x+c/x)/2$, first suppose
$1/2\le x\le5/2$. Normalize this stored dyadic exactly as $x=2^j d$,
where $1\le d<2$ and $j\in\{-1,0,1\}$. The reciprocal routine at
target $t$ returns $z$ with $|1-dz|\le2^{-t-1}$. Put $v=2^{-j}z$.
Because $d\ge1$ and $2^{-j}\le2$, we have
$|v-1/x|\le2^{-t}=\delta$. Compute the quotient and root update by
\[
 q_t=R_t(c_t v),\qquad x^+=R_t\bigl((x+q_t)/2\bigr).
\]
Here $c_t<4$ and $1/x\le2$, so the quotient has the uniform error
\[
 |q_t-c/x|
 \le\delta+c_t|v-1/x|+\frac{|c_t-c|}{x}
 \le7\delta.
\]
The exact Newton identity gives
\[
 |x^+-\sqrt c|
 \le\frac{(x-\sqrt c)^2}{2x}+\delta+\frac12|q_t-c/x|
 \le\frac{(x-\sqrt c)^2}{2x}+\frac92\delta.
\]
The fixed seed $x=3/2$ has $|x-\sqrt c|\le1/2$. At this seed the
first term is at most $1/12$, so one step with $t=6$ has error less
than $1/12+5/64<1/4$. Thereafter the invariant
$|x-\sqrt c|\le2^{-m-1}$, with $m\ge1$, gives
$3/4\le x\le9/4$ and hence supplies the denominator bound for the
next step. Since $1/(2x)\le1$ and $9/2<6$, the same capped doubling
schedule with five guard bits gives $|x-\sqrt c|\le2^{-p-1}$ at
target $p$.

These operations use a fixed number of scalar registers. In a reciprocal
step, $x\le3/2$ and $u\le c_t x\le3/2$ whenever $(*)$ is used. In a
root step, $x\le5/2$, $v\le2+\delta<3$, and $q_t<12$. A fixed number
of integer guard bits therefore suffices. At working precision $t$,
each retained significand has $t+O(1)$ bits, including the fixed
additional precision of the inner reciprocal, and each exact product
has $2t+O(1)$ bits. Theorem~\ref{thm:main} applies to the padded
significands. The streaming conventions give $O(t)$ cost for additions,
shifts, copying, truncation and head returns.

Each stage starts at the marked beginning of its normalized input, reads
only the required prefix, and returns over that prefix. If the finite
input ends first, the remaining fractional bits are zero. A root stage
forms $d$ once in $O(t)$ time; each stage of its inner reciprocal routine
likewise reads only its own shorter prefix of $d$. Thus no early stage
rescans an entire longer input.

For the capped doubling targets $m_j$ ending at $p$,
$\sum_j(m_j+5)=O(p)$, and monotonicity of $\lg$ gives
\[
 \sum_j B(m_j+5)
 \le(\lg(p+5))^{1-\kappa}\sum_j(m_j+5)
 =O(B(p)).
\]
The fixed warm-up adds constant cost. A reciprocal stage uses two
products, so the complete reciprocal routine at target $p$ costs
$O(B(p))$. A root stage of working precision $t$ uses a reciprocal
at target $t$ and one further product, and hence costs $O(B(t))$.
The outer precision sum is again $O(B(p))$. The scalar scan costs
sum to $O(p)$ at both levels.

The outer registers and the reused reciprocal and multiplier tapes form
one fixed finite tape set with fixed subroutine depth. Marking the ends
of each call's visited intervals on fixed tracks permits those intervals
to be cleared and their heads returned in a constant multiple of the
call's running time. No random-access simulation is used.

Take $P=2n+4$. At this target there are $O((\lg n)^2)$ outer and inner
stage calls. Their $O(\lg n)$-bit length and precision descriptors incur
a fixed polynomial in $\lg n$ of total work, hence $O(n)$. Together with
the initial normalization scan, this accounts for the fixed-tape
implementation. For division, the reciprocal routine gives
$|x-1/c|\le2^{-P}$. Form $a2^{-h}x$ by an exact significand product
and an exact shift of the binary point, retaining its fractional bits.
Since $a<2^n$,
\[
 \left|a2^{-h}x-\frac ab\right|
 \le a2^{-h}2^{-P}<2^{n-P}<\frac14.
\]
Its floor $q_0$ differs from $q=\lfloor a/b\rfloor$ by at most one.
Compute $r_0=a-bq_0$ exactly. If $r_0<0$, decrement $q_0$ and add $b$;
if $r_0\ge b$, increment $q_0$ and subtract $b$.
For square root, an approximation to $\sqrt c$ within $2^{-P}$ becomes
an approximation to $\sqrt a$ within $2^{e-P}<1/4$ after the exact
shift by $2^e$. Its floor $s_0$ differs from $\lfloor\sqrt a\rfloor$
by at most one. Decrement it if $s_0^2>a$, and increment it if
$(s_0+1)^2\le a$; otherwise retain it. The remainder and square tests
cover exact quotients and perfect squares. All final significands, trial
answers and exact test products have $O(n)$ bits; their computation costs
$O(B(n))$, and output padding is linear. Leading zeroes
affect neither normalization bound, since $n$ is the supplied input
length.
\end{proof}
